Put
\[
 G_V=\mathcal G[V].
\]
The premise is failure of ordinary identification on the positive finite substantive model class. Lemma \(\mathrm{lem{:}kivva\text{-}positive\text{-}ordinary\text{-}id\text{-}specialization}\), applied to \(G_V\) and the cell \((x,y)\), therefore gives finite recursive semi-Markovian substantive SCMs \(N_1,N_2\) compatible with \(G_V\) and a probability table \(p\) such that, for every \(v\in\mathcal X_V\),
\[
 P_{N_1}(V=v)=P_{N_2}(V=v)=p(v)>0,
\]
while
\[
 \psi_{x,y}(N_1)\ne\psi_{x,y}(N_2).
\]
This is the only step imported from the cited positive ordinary-identification result.

Fix any terminal selection kernel allowed by the statement and abbreviate
\[
 r_v=\pi_S(1\mid v_{\operatorname{pa}_{\mathcal G}(S)}),
 \qquad
 Z=\sum_{v'\in\mathcal X_V}r_{v'}p(v').
\]
Every \(r_v\) and every \(p(v)\) is strictly positive. The state space \(\mathcal X_V\) is finite by \(\mathrm{ass{:}finite\text{-}state}\), so
\[
 0<Z<\infty.
\]
Because no bidirected edge is incident to \(S\), using an exogenous selection coordinate independent of all substantive exogenous coordinates is compatible with \(\mathcal G\). Because \(S\) is terminal, adjoining its structural equation changes no substantive structural equation and creates no directed feedback. Hence each extension \(M_k\), \(k\in\{1,2\}\), is recursive, semi-Markovian, and compatible with \(\mathcal G\).

Conditional on \(V=v\), the independent common selection equation has success probability \(r_v\). Thus, for \(k\in\{1,2\}\), the product rule and the common population law give
\[
 P_{M_k}(V=v,S=1)
 =P_{M_k}(S=1\mid V=v)P_{M_k}(V=v)
 =r_vp(v).
\]
Summing this equality over the finite state space yields
\[
 P_{M_k}(S=1)=\sum_v r_vp(v)=Z>0.
\]
Division is therefore licensed, and
\[
 P_{M_k}(V=v\mid S=1)
 =\frac{r_vp(v)}{Z}
 =q_S(v).
\]
The right-hand side does not depend on \(k\). It is a probability table because
\[
 \sum_vq_S(v)=Z^{-1}\sum_vr_vp(v)=1,
\]
and every cell is strictly positive. Likewise marginalization of the common full law gives, for every \(t\in\mathcal X_{T_0}\),
\[
 P_{M_k}(T_0=t)
 =\sum_{v:v_{T_0}=t}P_{M_k}(V=v)
 =\sum_{v:v_{T_0}=t}p(v)
 =q_0(t).
\]
Consequently
\[
 \operatorname{Obs}_{T_0}(M_1)
 =\operatorname{Obs}_{T_0}(M_2)
 =(q_S,q_0).
\]

Finiteness and strict positivity of the selected table imply
\[
 \min_{v\in\mathcal X_V}q_S(v)>0.
\]
Fix \(\epsilon\) with \(0<\epsilon\leq\min_vq_S(v)\). The compatibility, positive selection normalizer, and selected-cell inequality established above put each \(M_k\) in \(\mathfrak M_\epsilon(\mathcal G)\) by Definition \(\mathrm{def{:}model\text{-}class}\). Hence the common observed pair lies in \(\mathcal C_\epsilon(\mathcal G,T_0)\) by Definition \(\mathrm{def{:}compatible\text{-}input\text{-}image}\), and Definition \(\mathrm{def{:}exact\text{-}observed\text{-}fiber}\) gives
\[
 M_1,M_2\in\mathfrak F_\epsilon(\mathcal G,T_0;q_S,q_0).
\]

For an intervention \(\operatorname{do}(X=x)\), terminality of \(S\) implies that the selection structural equation is downstream of every substantive variable and is not an input to any substantive equation. The substantive interventional distribution in \(M_k\) is therefore exactly that in \(N_k\). In particular,
\[
 \psi_{x,y}(M_k)=\psi_{x,y}(N_k),\qquad k\in\{1,2\},
\]
so the cited strict inequality yields
\[
 \psi_{x,y}(M_1)\ne\psi_{x,y}(M_2).
\]

It remains to verify, rather than assume, the response-interior assertion. Apply Lemma \(\mathrm{lem{:}finite\text{-}response\text{-}selection\text{-}admg\text{-}reduction}\) separately to \(M_1\) and \(M_2\). It gives finite independent response representations preserving their full factual laws, the \(T_0\)-margins, and the target intervention cells. In each representation merge response labels inducing the same deterministic maps and delete every zero-mass label. The resulting state space of each coordinate is exactly its positive support. Across the two representations there are only finitely many active coordinate atoms, and every one has strictly positive mass. Therefore
\[
 \eta=\min\{\text{active coordinate-atom masses in either reduced representation}\}>0.
\]
Every active atom in either representation has mass at least \(\eta\), while all equalities already proved are preserved by the reduction. This completes the graph-level positive selection lift.